\chapter{Verification layer}\label{ch:verify}

\section{The rule registry}
Every pass or fail verdict comes from \fpath{backend/app/verifier.py}. A rule is a plain function that receives the validated input and the \code{Design} and returns the fields of a \code{CheckResult}. The \code{@rule} decorator attaches the rule identifier, a title, an IS~456 clause, a severity (\code{error} or \code{warning}) and a flag \code{revisable} that states whether enlarging the section can fix a failure. Decorated rules are appended to one list, which is the single source for \code{run_rules}, for \code{GET /api/rules} and for the rule fingerprint. Table~\ref{tab:rules} lists the 19 rules.

\begin{longtable}{@{}L{3.55cm}L{4.35cm}L{3.3cm}cc@{}}
\caption{The 19 verifier rules, read from \code{RULE_META}. Sev.\ is the severity (E error, W warning); Rev.\ states whether resizing can fix a failure.}\label{tab:rules}\\
\toprule Rule identifier & Title & Clause & Sev. & Rev. \\ \midrule \endfirsthead
\toprule Rule identifier & Title & Clause & Sev. & Rev. \\ \midrule \endhead
\bottomrule \endlastfoot
\code{XCHK_MU} & Independent check of design moment & Table 18; \clause{22.2} & E & no \\
\code{XCHK_D} & Effective depth consistent with bar arrangement & \clause{26.4}; \clause{22.2} & E & no \\
\code{MU_LIM} & Moment within singly reinforced limit & \clause{38.1}; Annex G-1.1(a) & E & yes \\
\code{MOMENT_CAPACITY} & Moment of resistance of provided steel & Annex G-1.1(b) & E & yes \\
\code{XU_MAX} & Neutral axis within limiting depth & \clause{38.1} & E & yes \\
\code{MIN_TENSION_STEEL} & Minimum tension reinforcement & \clause{26.5.1.1(a)} & E & yes \\
\code{MAX_TENSION_STEEL} & Maximum tension reinforcement & \clause{26.5.1.1(b)} & E & yes \\
\code{TAU_V_MAX} & Maximum shear stress & \clause{40.2.3}; Table 20 & E & yes \\
\code{SHEAR_STIRRUPS} & Stirrup spacing provides shear strength & \clause{40.4(a)} & E & yes \\
\code{MIN_SHEAR_REINF} & Minimum shear reinforcement & \clause{26.5.1.6} & E & yes \\
\code{MAX_STIRRUP_SPACING} & Maximum stirrup spacing & \clause{26.5.1.5} & E & yes \\
\code{STIRRUP_PRACTICAL} & Stirrup spacing is practical to place & \clause{26.5.1.5}; good practice & E & yes \\
\code{DEFLECTION} & Span / effective depth (deflection control) & \clause{23.2.1}; Fig.~4 & E & yes \\
\code{COVER_EXPOSURE} & Nominal cover for exposure & \clause{26.4.2}; Table 16 & E & no \\
\code{GRADE_EXPOSURE} & Minimum concrete grade for exposure & Table 5 & E & no \\
\code{BAR_SPACING_MIN} & Minimum clear spacing between bars & \clause{26.3.2} & E & yes \\
\code{BAR_SPACING_MAX} & Maximum clear spacing of tension bars & \clause{26.3.3(b)}; Table 15 & E & yes \\
\code{SIDE_FACE_STEEL} & Side-face reinforcement for deep beams & \clause{26.5.1.3} & W & no \\
\code{STEEL_CONGESTION} & Tension steel percentage reasonable & \clause{26.5.1.1}; good practice & W & no \\
\end{longtable}

Of the 19 rules, 17 are error-level and two (\code{SIDE_FACE_STEEL}, \code{STEEL_CONGESTION}) are warnings, which are reported but never block a design. Four error-level rules are not revisable: the two cross-checks, and the two exposure rules, because a larger section cannot repair an under-specified cover or grade. Comparisons use a relative tolerance of $10^{-9}$ so that floating-point noise cannot change a verdict.

\section{Two independent cross-checks}
A verifier that reuses the calculator's own intermediate values would repeat the calculator's mistakes. Each rule therefore recomputes the quantity it checks from the raw input with its own formula; for example \code{MOMENT_CAPACITY} recomputes $M_{u,R}$ from the provided steel, and \code{XU_MAX} recomputes $x_u$. Two rules go further and exist only as cross-checks of the calculator.
\begin{itemize}
  \item \code{XCHK_MU} recomputes self-weight, factored load and design moment from the input alone, with its own constants, and requires agreement with the calculator's $M_u$ to within 0.5\,\%.
  \item \code{XCHK_D} requires the calculator's effective depth to equal $D - \bar y$ for the chosen bars to within 0.5~mm.
\end{itemize}
Because both are non-revisable, a calculator fault is reported as \code{needs_engineer} and is never repaired silently by resizing. The mutation tests of Chapter~\ref{ch:test} inject such faults to demonstrate that they are caught.

\section{The reviser loop}
The function \code{design_with_verification} runs the calculator and the 19 rules and then decides what to do (Figure~\ref{fig:loop}).
\begin{enumerate}
  \item If no error-level rule fails, the status is \code{designed}.
  \item If any failing error-level rule is not revisable, the loop stops with \code{needs_engineer}.
  \item If 8 revisions have already been made (\code{MAX_REVISER_ITERATIONS}), the loop stops with \code{needs_engineer}.
  \item Otherwise the reviser enlarges the section and the loop starts again.
\end{enumerate}
The reviser (\fpath{backend/app/reviser.py}) is deterministic. It increases the depth by 25~mm while the new depth satisfies $D + 25 \le 3b$ and $D + 25 \le 2000$~mm. When depth can no longer grow it increases the width by 25~mm while $b + 25 \le 1000$~mm. The enlarged section is validated again as a \code{BeamInput}. If neither step is possible the result is \code{needs_engineer} with the message that the section is at the supported size limits. Every iteration is recorded (section, failed rules, action) and is shown in the result as a revision history; a worked example is given in Appendix~\ref{app:example}.

\begin{figure}[H]
\centering
% Verify-and-revise loop (design_with_verification). Every node is deterministic code.
{\hyphenpenalty=10000\exhyphenpenalty=10000
\begin{tikzpicture}[x=1cm,y=1cm,
  dec/.style={det,diamond,aspect=3.2,inner sep=0pt,text width=3.3cm,font=\scriptsize},
  term/.style={neu,font=\small\bfseries}]
  \node[term] (in) at (4.5,0) {\texttt{BeamInput} (validated)};
  \node[det,text width=4.6cm] (calc) at (4.5,-1.5) {\texttt{design\_beam}: calculator, all numbers};
  \node[det,text width=4.6cm] (rules) at (4.5,-3.0) {\texttt{run\_rules}: 19 rules};
  \node[dec] (d1) at (4.5,-5.0) {any error-level rule failed?};
  \node[dec] (d2) at (4.5,-7.2) {all failing rules revisable?};
  \node[dec] (d3) at (4.5,-9.4) {8 revisions already made?};
  \node[det,text width=5.6cm] (rev) at (4.5,-11.4) {\texttt{revise}: depth +25 mm while $D+25\le3b$ and $\le2000$; else width +25 mm while $\le1000$};
  \node[dec] (d4) at (4.5,-13.6) {new section found?};
  \node[term,fill=detfill!55] (ok) at (11.0,-5.0) {designed};
  \node[term,fill=detfill!55,text width=3.3cm] (ne) at (11.0,-9.4) {needs\_engineer};

  \draw[flow] (in) -- (calc);
  \draw[flow] (calc) -- (rules);
  \draw[flow] (rules) -- (d1);
  \draw[flow] (d1) -- node[lbl,above] {no} (ok);
  \draw[flow] (d1) -- node[lbl,right] {yes} (d2);
  \draw[flow] (d2) -- node[lbl,right] {yes} (d3);
  \draw[flow] (d2.east) -- node[lbl,above,pos=0.25] {no} ++(2.4,0) |- ($(ne.north)+(0,0.0)$);
  \draw[flow] (d3.east) -- node[lbl,above] {yes} (ne.west);
  \draw[flow] (d3) -- node[lbl,right] {no} (rev);
  \draw[flow] (rev) -- (d4);
  \draw[flow] (d4.east) -- node[lbl,above,pos=0.3] {no} ++(3.3,0) -| (ne.south);
  \draw[flow] (d4.west) -- node[lbl,above,pos=0.3] {yes} ++(-2.6,0) |- (calc.west);
  \node[lbl,rotate=90] at (0.55,-7.5) {next iteration: calculator runs again};
\end{tikzpicture}}

\vspace{4pt}
{\small\detkey{} deterministic code; no model call occurs in this loop}

\caption{The verify-and-revise loop. All stages are deterministic code; no model is involved.}
\label{fig:loop}
\end{figure}

\section{Rule fingerprint}
\code{rules_fingerprint} is the first 12 hexadecimal characters of a SHA-256 hash over the string of \code{rule_id:clause} pairs joined by a vertical bar. It is stored with every run together with the code edition and the application, Python and Pydantic versions, so that a stored verdict can be tied to the rule set that produced it. In this version the fingerprint is \code{a9b7140d6bb9}. Adding or changing a rule changes the fingerprint.

\section{Section explorer}
The what-if explorer (\fpath{backend/app/explore.py}) passes neighbouring sections through exactly the same calculator and rules. For a base section it takes widths $b-50$, $b$ and $b+50$~mm (clamped to 150 to 1000~mm and de-duplicated) and depths from $D-100$ to $D+300$~mm in 25~mm steps, which gives at most $3 \times 17 = 51$ candidate sections; candidates that fail \code{BeamInput} validation are dropped. Rows are sorted by concrete area and then by steel area, and the lightest passing section is reported. The explorer uses no model and writes nothing to the audit log. It is not a cost optimisation.
